\par\addvspace{7pt}\Needspace{4\baselineskip}\noindent\textbf{稳定匹配的共同方法}\par
下面补充若干题共用的延迟接受证明。男女各有 $n$ 人，每人对另一侧的所有人作严格完整排序；“更喜欢”始终指这个排序。匹配 $M$ 的\textbf{阻塞对}是没有互相匹配、却各自更喜欢对方而非 $M$ 中伴侣的一男一女。没有阻塞对的完美匹配称为稳定匹配。

\begin{lemma}\label{lem:6042-stable}
男方提出申请的延迟接受算法终止于稳定完美匹配；每位男方得到他在所有稳定匹配中最喜欢的可行伴侣，每位女方得到她在所有稳定匹配中最不喜欢的可行伴侣。交换两侧角色得到另一极端。稳定匹配唯一，当且仅当两个极端匹配相同。
\end{lemma}
\begin{proof}
原课的按日重访版本是：每位男方每天到自己名单中尚未被划去的第一位女方处，包括仍被保留的男方也每日重访；女方在当日申请者中保留最喜欢的一位、拒绝其余人，被拒者永久划去她的名字。某一日无人被拒即停止。另一个等价的逐次版本只令尚未被女方保留的男方提出新的申请；女方在新申请者及此前保留者中择优。两版中女方一旦有申请者，此后一直保留一人，而且保留者只会变得更合意；男方只会保持原对象或向名单中较后的人移动。下文在线的按日不变量按第一个版本理解。

按日重访版中每个发生拒绝的日子至少永久划去一个名字，总共至多 $n^2$ 次划去，所以有限日内必到达无人被拒的日子并停止。逐次版中每对男女至多发生一次新申请，故至多 $n^2$ 次新申请后终止。若某男方的名单已空，他被全部 $n$ 位女方拒绝，全部女方已经保留各不相同的男方；这需要除他外另有 $n$ 位男方，与总数 $n$ 矛盾。因此没有男方名单耗尽，终止时所有人匹配。若男方 $b$ 比起终局伴侣更喜欢 $g$，他此前必已向 $g$ 申请并被拒绝；$g$ 的终局伴侣比 $b$ 更合意，故 $(b,g)$ 不是阻塞对。这证明稳定性。

关键是证明\textbf{任何曾被拒绝的配对都不可能出现在任何稳定匹配中}。反设首次发生这种拒绝时，女方 $g$ 拒绝 $b$ 而保留 $c$，且存在稳定匹配 $S$ 把 $b$ 配给 $g$。$g$ 更喜欢 $c$。若 $c$ 更喜欢 $S(c)$ 而非 $g$，则 $c$ 在申请 $g$ 前已被 $S(c)$ 拒绝，而 $(c,S(c))$ 是稳定匹配 $S$ 中的配对，违反“首次”的选择。因此 $c$ 更喜欢 $g$ 而非 $S(c)$；于是 $(c,g)$ 阻塞 $S$，矛盾。证明成立。

每个男方排在终局伴侣之前的人都已拒绝他，故终局伴侣就是他最喜欢的稳定可行伴侣，且不依赖申请处理次序。再比较任意稳定匹配 $S$ 与男方最优匹配 $M$。若某女方 $g$ 更喜欢 $M(g)=b$ 而非 $S(g)$，那么 $b$ 因男方最优性也更喜欢 $g=M(b)$ 而非 $S(b)$，产生阻塞对。因此 $g$ 在任意 $S$ 中的伴侣至少与 $M(g)$ 同样合意，证明女方最差性。角色交换给出女方最优、男方最差匹配。每位男方在任意稳定匹配中的伴侣都夹在两个极端之间；两个极端相同则所有配对被迫相同，反向显然。
\end{proof}

本节讨论“保持不变量”时，意思是\textbf{某谓词在一个状态为真，则下一轮仍为真}，并不要求它初始就为真。一个在所有可达状态恒假的谓词也满足这个蕴含。按日同时申请的版本中，有些日子无人被拒绝；因此名单长度通常只弱单调，不一定每一日都严格变化。
